\section{Discussion}\label{sec:discussion}

\subsection{More rows are not more information}
The central finding of this study is a distinction between sample size and independent
information. The corpus contains 118 rows but only 40 independent experimental contexts,
and 64\% of the variance in capacitance lies between those contexts rather than
within them. Bayesian-network resampling multiplies rows while holding the number of contexts
fixed. Since the generator's joint distribution is estimated from the training papers alone,
every synthetic observation is a recombination of dependencies those papers already exhibit.
When the test partition is a \emph{different} publication --- with its own protocol, its own
mass-normalisation convention and, as Sec.~\ref{sec:domainshift} shows, frequently its own
unique descriptor settings --- no amount of resampling can supply the missing information.
The augmentation null result and the fidelity success are therefore two aspects of the same
fact rather than a contradiction.

This also delimits when the motivating FCDI-style strategy should be expected to work. Where
observations are exchangeable draws from one experimental campaign, resampling the joint
distribution is a reasonable regulariser. Where they are clustered by study and evaluation is
across studies, it addresses the wrong bottleneck.

\subsection{Four purposes that must not be mixed}
The analysis maintains four separate roles, and much of the interpretive risk in this problem
comes from conflating them. \emph{Predictive assessment} uses paper-grouped
cross-validation. \emph{Descriptor interpretation} uses the final real-only model fitted to
all observations. \emph{Robustness testing} uses Bayesian augmentation and SHAP-ranking
stability. \emph{Physics-model design} uses descriptors supported by all three. In
particular, the in-sample fit of the interpretation model is never quoted as performance, and
the stability of the SHAP ranking is never taken as evidence that the model predicts well ---
it does not.

\subsection{Attribution, association and causation}
SHAP quantifies how a fitted model distributes credit among its inputs. Within this corpus
that attribution is stable and reproducible, which makes it a defensible basis for
prioritising descriptors. It is not evidence of causation, and Sec.~\ref{sec:shap} gives a
quantitative reason for caution specific to this dataset: several highly ranked descriptors
are paper-level constants (ICC$(1) \approx 1$), so their attribution is statistically
entangled with publication identity. The interlayer-spacing trend, which runs opposite to the
naive physical expectation, is the clearest symptom. The appropriate conclusion is that these
descriptors organise the corpus and the model uses them, not that changing them would change
measured capacitance.

\subsection{What the physics analysis changes}
Rejecting a classical PINN is a substantive result rather than a negative one. It redirects
effort from constructing a differential residual the data cannot support towards two things
the data do support: enforcing the verified closure identities of
Eqs.~\eqref{eq:areal}--\eqref{eq:vol} as internal-consistency constraints, and extending
the extraction schema to recover the potential window and discharge time. The second is the
higher-value action. It converts the physics term from a unit-consistency relation applicable
to a minority of rows into the measurement equation itself, and it simultaneously addresses
the heterogeneity diagnosis, because knowing $\Delta V$ and the testing mode removes two known
and currently unmodelled sources of between-study scatter.

\subsection{Implications for reporting practice}
The rows discarded during cleaning are themselves a finding. Capacitance reported only as a
range or an inequality, current densities given on incompatible bases, mass loadings without
an area, and flake sizes given qualitatively together removed a substantial fraction of
otherwise usable records. Standardised reporting of the potential window, testing mode,
discharge time and normalisation basis would improve the value of the literature for
data-driven analysis more than any modelling refinement available on the present corpus.
